% discussion.tex (v15, 2026-10-21): no subheadings; one argument in six paragraphs: what the anchor does, why the Alaska gain is small, what the product comparison shows and why, how this relates to previous work, what the maps and the label-number results mean for a new region, limitations.
\section{Discussion}\label{sec:discussion}

The main result is that the position at which physical knowledge enters the model determines its accuracy within regions and its transferability between regions. In the physics-residual formulation (equation~(\ref{eq:residual})) the Stefan term represents the response to the thaw-season climate of each year as proportional to the square root of TDD, and the learner corrects the departures from one regional coefficient that soils and terrain impose. This division of labour explains the main patterns of the results. The gain over the Stefan model appeared only when terrain and soil covariates were available, which is consistent with these variables setting the effective coefficient of each site; it increased with the error of the Stefan model; and terrain variables accounted for \nshapLeResTerrain\% of the SHAP attribution of the residual model in the Lena Delta and soil variables for \nshapCaResSoil\% in Canada. These patterns indicate that the residual learner captures the part of the thaw depth that one Stefan coefficient leaves unexplained. Direct ML, which must learn the climate response from the labels, had a lower site-year RMSE with thousands of labels, but with 25 labels it was less accurate than the Stefan model in Alaska and Canada and, when the Lena Delta was withheld, its error exceeded that of the Stefan model by \ndrefLeregDzabs{} cm. The in-region RMSE of the physics-residual model exceeded that of direct ML by at most \ndrefRoDzblkmax{} cm, while the model kept the transferability of its anchor, which matters where labels are scarce.

In Alaska the residual correction produced no gain, for two reasons. The recalibrated Stefan model already fits the regional level and the year-to-year variation of the dense Alaska labels, so that little systematic residual is left for the covariates to explain, and the error that remains is mostly the spread between probe points inside one 1 km cell (pooled standard deviation \ncyWithinSdAk{} cm), which no covariate at 1 km resolution can resolve. The absence of a gain in every tercile (Fig.~\ref{fig:conditions}b) and the same result on cell-year means (Supplementary Table~S5) are consistent with these two reasons and do not indicate a failure of the learner. The learner is therefore not always needed: the spatial-block comparison between the physics-residual model and the recalibrated Stefan model, computed with the labels of a region, shows whether that region has systematic residuals that covariates can explain (a resolved negative $\Delta$RMSE, as in the Lena Delta and Canada) or not (Alaska), and the learner is worth adding only in the first case.

On the same sites and years, the accuracy of the ML products depended on the distance to their training sites, which validations within their own site sets do not assess. Where their training sites are close to our label sites (Alaska, and Canada for Liu et al., whose training data include the earlier release of the same observations), the ML products were as accurate as our model, and the comparison there is not independent. Where they have no nearby training sites (the Lena Delta, median distance \ndistweiLemed{} km), they predicted between about 40 and 70 cm regardless of the observed ALT, overestimated ALT in the Lena Delta by about 20 cm, most strongly at the thin active layers of the northern delta, and were less accurate than a physics-residual model trained without any Lena Delta label. This pattern matches the error increase of direct ML under region hold-out; we attribute it to products that learn the climate response at training sites elsewhere, although their predictors are yearly. In the physics-residual model the response to the TDD of the stated year is carried by the Stefan term, which does not depend on the training sites, so this error increase is small. The process-based ESA CCI product, which is not trained on ALT observations, was less accurate than our model in all three regions but, unlike the ML products, was not systematically biased in the Lena Delta (Fig.~\ref{fig:products}c), which is consistent with its physical basis.

These results extend previous comparisons. Gautam et al.\cite{Gautam2025} found at 68 Alaska CALM sites of one year that a random forest lost most of its accuracy on held-out sites whereas the Stefan model retained it; this ordering holds here over site-years of 2010--2024 in three regions and under region hold-out, and the two approaches are complementary because a learner that corrects the Stefan model retains its transferability. Studies of physics-guided learning in lakes and rivers reported that process knowledge helps mainly under data scarcity and extrapolation\cite{Read2019,Jia2021,Willard2022}; the label-number and region hold-out experiments quantify this for ALT. The finding that random cross-validation underestimates the spatial-block error of direct ML by \noptimLeDz--\noptimCaDz{} cm agrees with the literature on spatial validation\cite{Roberts2017,Ploton2020,MeyerPebesma2022}; together with the region hold-out results, this indicates that errors from random or site-based splits do not represent the accuracy of a product in a region without training observations.

The yearly maps apply the physics-residual model to the three regions with an area of applicability and prediction intervals as conditions of use. Each map uses the ERA5-Land climate of one year with the TDD accumulated to 1~September and can therefore be compared with observations or products of the same year. The regional mean anomalies of the maps are the modelled response to the thaw-season climate (correlation with the regional TDD \ntddanomAk{}--\ntddanomCa{}). They can be checked against observations only at sites measured year after year (Supplementary Fig.~S9). The linear trends of the maps are the modelled response to the climate of 2011--2024, not observed trends. Without local observations, the Stefan model with a coefficient from other regions or the physics-residual model trained on other regions can serve as the starting point; neither showed the large error increase of direct ML in the Lena Delta, although in Canada direct ML trained on other regions was the most accurate. The effect of 25 local labels differed between regions: they reduced the error of the physics-residual model from \ntlrmseLeRszero{} to \ntlrmseLeRstwfive{} cm in the Lena Delta but increased it from \ntlrmseCaRszero{} to \ntlrmseCaRstwfive{} cm in Canada, which we attribute to the difference in thaw depth between the two Canadian data sources. The spatial-block comparison within the new region and the area of applicability of the resulting map distinguish these cases.

The study has limitations. Half of the Alaska labels come from 2014, which prevented a resolved estimate of the effect of the year-specific input there, and the Lena Delta labels come from one data set in \nnblkLe{} blocks, which widens the intervals. The labels combine probe transects and CALM grid means from three programmes, and the model error differs between them (Supplementary Table~S7), most clearly at the \nsrcCaCALMN{} Canadian CALM site-years, whose mean ALT of \nsrcCaCALMObs{} cm is underestimated by \nsrcCaCALMBias{} cm; differences between measurement protocols and site selection are therefore part of the reported errors. The labels are observation sites rather than independent 1~km cells, which the spatial-block design and the cell-year scores account for. The climate input is monthly and the soil covariates are static. The Kudryavtsev model, which includes snow and moisture terms, was more accurate than the recalibrated Stefan model in Canada, and most of the difference between the two anchors there was corrected by the residual learner; a moisture-aware anchor might further reduce the error that remains after this correction. Active layers thicker than about 3~m lie outside the training range, and the comparison of maps with observations in one year rests on few cells in the Lena Delta and Canada. Because every number is tied to a stated site, year and validation design, the protocol can be applied to new products and regions, and the maps extended year by year.
